\ifdefined\CORRECTION
\documentclass[correction]{td}
\else
\documentclass{td}
\fi

\inputexercisepath[src/corrections/TP]{src/exercises}

\usepackage[T1]{fontenc}
\usepackage[hidelinks]{hyperref}
\usepackage{amsmath,amssymb, mathrsfs, stmaryrd}
\newcommand\eqdef{\stackrel{\text{def}}{=}}

\usepackage{algos}
\usepackage{statemachines}
\usepackage{trees}

\codeUE{XLG4IU020}
\intituleUE{Langages et Automates}
\cursus{L2 Informatique}

\author{Matthieu \textsc{Perrin}}

\logo{src/img/logoUN.png}
\institution{Nantes Université}

\typeTP[4]
\title{Analyse statique de programmes}

\hypersetup{
  pdftitle  = {Langages et Automates - TP 4 - Analyse statique de programmes},
  pdfauthor = {Matthieu Perrin},
  pdfsubject  = {TP de L2 en Langages et Automates},
  pdfkeywords = {analyse statique, typage, initialisation des variables}
}

\begin{document}

\maketitle

Dans ce TP, nous nous intéressons à \emph{l'analyse statique de programmes}.
Contrairement à l'interprétation, qui exécute un programme sur des valeurs concrètes,
l'analyse statique consiste à raisonner \emph{sans exécuter le programme},
afin de détecter des erreurs potentielles ou des propriétés garanties.

\begin{exercice}[Mise en place]

  \begin{question}
  \item Clonez le dépôt du TP :
    \begin{lstlisting}[gobble=6, language=bash]
      git clone https://github.com/LangagesEtAutomates/lea-tp5-analyser
    \end{lstlisting}

    La compilation et l'exécution de l'interpréteur sont similaires aux TPs précédents.
  \end{question}

  \paragraph{Analyseur statique fourni}
  
  La classe \lstinline{lea.Analyser} implémente une analyse statique de détection d'utilisation de variables avant initialisation ou de variables non déclarées.
  L'analyse parcourt l'AST et maintient un contexte de type \lstinline{Context} qui contient deux ensembles d'identifiants :
  \begin{itemize}
  \item \lstinline{declared} :
    l'ensemble des variables déclarées au début du programme et inchangé au cours de l'analyse,
    qui permet de vérifier que toute variable a été déclarée avant d'être utilisée ;
  \item \lstinline{written} :
    l'ensemble des variables \emph{certainement initialisées},
    qui permet de vérifier que toute variable a été initialisée avant d'être lue.
  \end{itemize}
  
  L'analyse est réalisée par une fonction récursive qui transforme un contexte en un
  nouveau contexte enrichi par les effets de l'instruction analysée. 
  \begin{description}
  \item[L'affectation :]
    L'analyse d'une affectation (méthode \lstinline{Context analyse(Assignment assignment, Context context)})
    commence par analyser que la variable à affecter (\lstinline{assignment.lhs}) est bien déclarée,
    et que l'expression affectée (\lstinline{assignment.rhs}) n'utilise pas de variables non déclarées ou initialisées.
    Si tout est correct, elle retourne un contexte où la variable est ajoutée à l'ensemble \lstinline{context.written}.
    
    Par exemple, si $\text{\lstinline{context.written}} = \{x\}$ avant l'affectation \lstinline{y <- 2;},
    $\text{\lstinline{context.written}} = \{x, y\}$ dans le contexte retourné après l'affectation.

  \item[La séquence :]
    Dans une suite d'instructions (méthode \lstinline{Context analyse(Sequence sequence, Context context)}),
    le contexte produit par une commande devient le contexte d'entrée de la commande suivante
    pour modéliser la propagation linéaire de l'information.

    Par exemple, si $\text{\lstinline{context.written}} = \emptyset$ avant la séquence \lstinline{x <- 1; y <- 2;},
    alors $\text{\lstinline{context.written}} = \{x\}$ entre les deux affectations, et $\text{\lstinline{context.written}} = \{x, y\}$
    dans le contexte retourné à la fin de la séquence.
  \end{description}  

  Nous nous intéressons maintenant à l'analyse des structures conditionnelles. 
  L'analyse est dite \emph{pessimiste} : elle ne valide un programme que si l'initialisation est garantie pour \emph{tous} les chemins d'exécution possibles.
  Ainsi, à la jonction de deux branches d'une conditionnelle, l'ensemble des variables considérées comme initialisées est l'\emph{intersection} des variables
  initialisées dans chaque branche.  

  \pagebreak
  \begin{question}
  \item Indiquez quelles erreurs sont signalées par l'analyseur dans les programmes suivants, et pourquoi.
    
    \begin{minipage}{.3\textwidth}
    \begin{lstlisting}
      algorithme
      variables
        x;
      début
        x <- x + 1;
      fin
    \end{lstlisting}
    \end{minipage}%
    \begin{minipage}{.3\textwidth}
    \begin{lstlisting}
      algorithme
      variables
      début
        x <- 0;
        x <- x + 1;
      fin
    \end{lstlisting}
    \end{minipage}%
    \begin{minipage}{.3\textwidth}
    \begin{lstlisting}
      algorithme
      variables
        x;
      début
        si 1+1=2 alors
          x <- 0;
        sinon
          écrire(x);
        fin si
        x <- x+1;
      fin
    \end{lstlisting}
    \end{minipage}%
  \end{question}
  \begin{correction}
    \begin{minipage}{.3\textwidth}
    \begin{lstlisting}
      algorithme
      variables
        x;
      début
      // x pas initialisée
      x <- x + 1;
      fin
    \end{lstlisting}
    \end{minipage}%
    \begin{minipage}{.3\textwidth}
    \begin{lstlisting}
      algorithme
      variables
      début
        // x pas déclarée
        x <- 0;
        // x pas déclarée
        x <- x + 1;
      fin
    \end{lstlisting}
    \end{minipage}%
    \begin{minipage}{.3\textwidth}
    \begin{lstlisting}
      algorithme
      variables
        x;
      début
        si 1+1=2 alors
          x <- 0;
        sinon
          // x pas initialisée
          écrire(x);   
        fin si
        // x pas initialisée
        x <- x+1;
      fin
    \end{lstlisting}
    \end{minipage}%
  \end{correction}

  \begin{question}
  \item Étudiez la méthode \lstinline{analyse(If i, Context context)}. 
    Pour chacun des exemples suivants, identifiez la valeur de \lstinline{written} après chaque instruction, 
    et justifiez pourquoi l'analyse considère qu'une variable
    est ou n'est pas initialisée après la conditionnelle.

    \begin{minipage}{.3\textwidth}
      \begin{lstlisting}[gobble=8]
        algorithme
        variables
          x;
        début
          si vrai alors
            x <- 0;
          sinon
            écrire(1);
          fin si
          écrire(x);
        fin
      \end{lstlisting}
    \end{minipage}%
    \begin{minipage}{.3\textwidth}
      \begin{lstlisting}[gobble=8]
        algorithme
        variables
          x;
        début
          si vrai alors
            x <- 0;
          sinon
            x <- 1;
          fin si
          écrire(x);
        fin
      \end{lstlisting}
    \end{minipage}%
    \begin{minipage}{.3\textwidth}
      \begin{lstlisting}[gobble=8]
        algorithme
        variables
          x; y;
        début
          si vrai alors
            x <- 0;
          sinon
            y <- 1;
          fin si
          écrire(x);
          écrire(y);
        fin
      \end{lstlisting}
    \end{minipage}%
  \end{question}
  \begin{correction}
    \begin{minipage}{.3\textwidth}
      \begin{lstlisting}[gobble=8]
        algorithme
        variables
          x;
        début           // {}
          si vrai alors
            x <- 0;     // {x}
          sinon
            écrire(1);  // {}
          fin si        // {}
          écrire(x);    // erreur
        fin
      \end{lstlisting}
    \end{minipage}%
    \begin{minipage}{.3\textwidth}
      \begin{lstlisting}[gobble=8]
        algorithme
        variables
          x;
        début            // {}
          si vrai alors
            x <- 0;      // {x}
          sinon
            x <- 1;      // {x}
          fin si         // {x}
          écrire(x);     // {x}
        fin
      \end{lstlisting}
    \end{minipage}%
    \begin{minipage}{.3\textwidth}
      \begin{lstlisting}[gobble=8]
        algorithme
        variables
          x; y;
        début            // {}
          si vrai alors
            x <- 0;      // {x}
          sinon
            y <- 1;      // {y}
          fin si         // {}
          écrire(x);     // erreur
          écrire(y);     // erreur
        fin
      \end{lstlisting}
    \end{minipage}%
  \end{correction}
\end{exercice}


\begin{exercice}[Détection de code mort]

  Le but de cet exercice est d'ajouter à notre langage
  une instruction d'\emph{interruption de boucle} \lstinline|interrompre|,
  analogue à \lstinline|break| en C/Java.
  Cette instruction ne peut apparaître qu'à l'intérieur d'une boucle.
  Lorsqu'elle est exécutée, elle termine \emph{immédiatement}
  l'itération courante et provoque la sortie de la boucle.

  \begin{question}
  \item Ajoutez l'instruction \lstinline|interrompre| au langage.
    \begin{itemize}
    \item Étendez la hiérarchie de types avec un nouveau type d'instructions \lstinline|Break|.
    \item Modifiez le lexer JFlex afin de reconnaître le nouveau mot-clé \lstinline|interrompre|.
    \item Modifiez le parseur CUP afin d'ajouter une règle de grammaire pour la nouvelle commande.
    \item Modifiez l'interpréteur afin de prendre en compte l'instruction \lstinline|interrompre|.

      \emph{Indice :} une solution simple consiste à lever une exception \lstinline|BreakException| (à créer)
      lors de l'interprétation d'une instruction \lstinline|Break|, et à la rattraper dans
      l'interprétation de la boucle pour en sortir proprement. 
    \end{itemize}
  \end{question}
  \begin{correction}
    \begin{lstlisting}[gobble=6]
      // Node.java
      public record Break() implements Instruction { }

      // Lexer.flex
      "interrompre" { return mark(Terminal.INTERROMPRE); }

      // Parser.cup
      commande ::= ...
        | INTERROMPRE:i {: RESULT = mark(new Break(), i, i); :};
      
      // Interpreter.java
      public void execute(Program program) {
	try {
	  try {
	    interpret(program.body());
	  } catch (BreakException e) {
	    error(e.node, "Interrompre ne peut pas être en dehors d'une boucle");
	  }
	} catch (PanicException e) {}
      }

      private void interpret(Instruction instruction) throws BreakException, PanicException {
	switch(instruction) {
          ...
	  case Break b 		-> throw new BreakException(b);
	}
      }
      
      private void interpret(While w) throws BreakException, PanicException {
	try {
	  while(evalAsBool(w.cond())) {
	    interpret(w.body());
	  }
	}catch(BreakException e) {}
      }
      // Idem pour la boucle pour...
    \end{lstlisting}
  \end{correction}

  Les instructions situées après un \lstinline|interrompre| 
  dans une même séquence d'instructions ne peuvent jamais être exécutées :
  elles constituent donc du \emph{code mort}.
  Par exemple, le code suivant dit signaler du code mort pour l'instruction \lstinline|écrire(999);| :
  \begin{lstlisting}
    algorithme
    variables
      x;
    début
      tant que x < 10 faire
        écrire(x);
        interrompre;
        écrire(999);   // code mort
      fin tant que
      écrire(0);       // atteignable
    fin
  \end{lstlisting}

  \begin{question}
  \item Quelles instructions des programmes suivants constituent-elles du \emph{code mort} ?

    \begin{minipage}{.33\textwidth}
      \begin{lstlisting}[gobble=8]
        algorithme
        variables
          x;
        début
          tant que x < 10 faire
            interrompre;
            écrire(1);
          fin tant que
          écrire(2);
        fin
      \end{lstlisting}
    \end{minipage}%
    \begin{minipage}{.33\textwidth}
      \begin{lstlisting}[gobble=8]
	algorithme
	  variables
	    x;
	  début
	    x <- 0;
	    tant que x < 10 faire
	      si x < 5 alors
		interrompre;
		écrire(1);
	      sinon
		écrire(2);
	      fin si
	      écrire(3);
	    fin tant que
	écrire(4);
	fin
      \end{lstlisting}
    \end{minipage}%
    \begin{minipage}{.33\textwidth}
      \begin{lstlisting}[gobble=8]
	algorithme
	variables
	  x;
	début
	  x <- 0;
	  tant que x < 10 faire
	    si x < 5 alors
	      interrompre;
	    sinon
	      interrompre;
	    fin si
	    écrire(1);
	  fin tant que
	  écrire(2);
	fin
      \end{lstlisting}
    \end{minipage}
  \end{question}
  \begin{correction}
        \begin{minipage}{.32\textwidth}
      \begin{lstlisting}[gobble=8]
        algorithme
        variables
          x;
        début
          tant que x < 10 faire
            interrompre;
            écrire(1); // code mort
          fin tant que
          écrire(2);  // accessible
        fin
      \end{lstlisting}
    \end{minipage}%
    \begin{minipage}{.36\textwidth}
      \begin{lstlisting}[gobble=8]
	algorithme
	  variables
	    x;
	  début
	    x <- 0;
	    tant que x < 10 faire
	      si x < 5 alors
		interrompre;
		écrire(1); // code mort
	      sinon
		écrire(2);
	      fin si
	      écrire(3);  // accessible
	    fin tant que
	  écrire(4);      // accessible
	fin
      \end{lstlisting}
    \end{minipage}%
    \begin{minipage}{.32\textwidth}
      \begin{lstlisting}[gobble=8]
	algorithme
	variables
	  x;
	début
	  x <- 0;
	  tant que x < 10 faire
	    si x < 5 alors
	      interrompre;
	    sinon
	      interrompre;
	    fin si
	    écrire(1); // code mort
	  fin tant que
	  écrire(2);  // accessible
	fin
      \end{lstlisting}
    \end{minipage}
  \end{correction}
    
  On souhaite enrichir l'analyse statique afin de signaler le code mort. Pour cela, on enrichit le contexte d'un booléen
  \lstinline|Analyser.Context.accessible|, initialement \lstinline|true|,
  est mis à \lstinline|false| lorsqu'on analyse \lstinline|interrompre|.
  À la sortie d'une structure de contrôle (boucle ou conditionnelle), le booléen doit être mis à jour pour détecter si
  les instructions suivantes sont \emph{forcément} du code mort, c'est-à-dire quelle que soit la branche exécutée. 

  \begin{question}
  \item Adaptez l'analyse statique afin de détecter lorsqu'une instruction est \emph{forcément}
    située après une instruction \lstinline|interrompre|.
    Lorsqu'une instruction \lstinline|commande| est analysée avec \lstinline|accessible = false|,
    une erreur doit être signalée en utilisant : 
    \lstinline|error(commande, "Code mort", context);|
  \end{question}
  \begin{correction}
    \begin{lstlisting}[gobble=6]
      public final class Analyser {
        ...
        private Context analyse(Node node, Context context) {
	  return switch(node) {
            ...
	    case Break b -> context.withAccessible(false);
            ...
	  };
        }
        ...
        private Context analyse(Sequence sequence, Context context) {
	  for(var commande : sequence.commands()) {
	    if(!context.accessible) error(commande, "Code mort", context);
	    context = analyse(commande, context);
	  }
	  return context;
        }
        ...
        static final class Context {
          ...
	  final boolean accessible;
	  public Context(Program program) {
	    declared = Map.copyOf(program.declared());
	    written = Set.of();
	    accessible=true;
	  }
          ...
	  public Context withAccessible(boolean accessible) { 
	    return new Context(declared, written, accessible); 
	  }
          ...
	  public Context merge(Context other) {
	    var writ = new HashSet<>(written);
	    writ.retainAll(other.written);
	    return new Context(declared, writ, accessible || other.accessible);
	  }
        }
      }
    \end{lstlisting}
  \end{correction}

\end{exercice}

\end{document}
